\appendix
\section{Data tables}\label{app:data}
The 51 records of spoilage flora are listed in Table~\ref{tab:dsp}. The 99 \textit{L.~monocytogenes} records, with their set (fitting or held out) and source, are listed in \ListeriaRef, and the complete compilation of \DatNFull{} records is in the repository.
\begin{table}[!htbp]\centering\scriptsize
\caption{ComBase records of aerobic spoilage flora in pasteurised milk ($a_w$ 0.997, pH 6.6; Tasmanian Institute of Agriculture). $\mu_{\max}$: maximum growth rate in decimal-logarithm units (log$_{10}$\,CFU\,mL$^{-1}$ per hour). Data from ComBase \cite{combase}}\label{tab:dsp}
\setlength{\tabcolsep}{4pt}
\begin{tabular}{@{}cc|cc|cc@{}}\toprule\rowcolor{dghead}
$T$/\textdegree C & $\mu_{\max}$/h$^{-1}$ & $T$/\textdegree C & $\mu_{\max}$/h$^{-1}$ & $T$/\textdegree C & $\mu_{\max}$/h$^{-1}$\\\midrule
2.7 & 0.0171 & 6.8 & 0.0242 & 10.7 & 0.0866\\
3.4 & 0.0139 & 7.3 & 0.0613 & 11.1 & 0.0910\\
3.4 & 0.0264 & 7.4 & 0.0551 & 11.5 & 0.0852\\
4.2 & 0.0334 & 8.1 & 0.0605 & 11.5 & 0.1040\\
4.2 & 0.0313 & 8.2 & 0.0653 & 11.6 & 0.0948\\
4.3 & 0.0426 & 8.7 & 0.0615 & 12.1 & 0.0961\\
4.6 & 0.0391 & 8.8 & 0.0706 & 12.5 & 0.1070\\
4.9 & 0.0360 & 8.8 & 0.0676 & 12.6 & 0.0750\\
5.0 & 0.0396 & 9.3 & 0.0623 & 12.6 & 0.1270\\
5.2 & 0.0489 & 9.4 & 0.0826 & 12.7 & 0.1040\\
5.7 & 0.0492 & 9.6 & 0.0749 & 13.2 & 0.1330\\
5.8 & 0.0520 & 9.6 & 0.0658 & 13.3 & 0.1110\\
5.9 & 0.0431 & 9.9 & 0.0859 & 13.8 & 0.1230\\
5.9 & 0.0519 & 10.1 & 0.0804 & 13.9 & 0.0910\\
6.3 & 0.0512 & 10.2 & 0.0715 & 14.2 & 0.1310\\
6.6 & 0.0486 & 10.5 & 0.0813 & 14.8 & 0.1500\\
6.6 & 0.0562 & 10.7 & 0.0740 & 14.9 & 0.1200\\
\bottomrule\end{tabular}\end{table}



\section{Proofs}\label{app:proofs}
\textbf{Proposition~\ref{prop:clock}.} Where $\mu>0$, $\tau(t)$ is strictly increasing and $\dd\tau=\mu\dd t$. Dividing Eqs.~\eqref{eq:bar}--\eqref{eq:barq} by $\mu$ gives $\dd y/\dd\tau=\frac{q}{1+q}(1-\mathrm{e}^{y-y_{\max}})$ and $\dd q/\dd\tau=q$, an autonomous system with no reference to $T$. Its solution $(Y,Q)$ from $(y_0,q_0)$ is unique (the right-hand side is Lipschitz), so $y(t)=Y(\tau(t))$ and $q(t)=q_0\mathrm{e}^{\tau(t)}$. Since $Y$ is increasing while $y<y_{\max}$, $N$ reaches $N_L$ exactly when $\tau$ reaches $\tau^*=Y^{-1}(\ln N_L)$, which depends on $(h_0,y_0,y_{\max},N_L)$ only. Intervals with $\mu=0$ ($T\le T_{\min}$) leave $\tau$, $y$ and $q$ unchanged. \hfill$\square$

\textbf{Proposition~\ref{prop:csl}.} From Eq.~\eqref{eq:sl}, $1/\SL(T)=\mu(T)/K$, and the continuous limit of Eq.~\eqref{eq:csl} is $\int_0^t\mu(T(s))\dd s/K=\tau(t)/K$. \hfill$\square$

\textbf{Corollary~\ref{cor:dh}.} Write $T-T_{\min}=(T_r-T_{\min})+(T-T_r)$, square, multiply by $b^2/K$ and integrate. \hfill$\square$

\textbf{Proposition~\ref{prop:calo}.} With $P=Y_QV\,\dd N/\dd t$ and $\dd N/\dd t=N\,\dd y/\dd t$, Eq.~\eqref{eq:bar} gives $P=Y_QV\mu(T)N\frac{q}{1+q}(1-N/N_{\max})$. By Proposition~\ref{prop:clock}, $N$ and $q$ are functions of $\tau$ alone, so $P/\mu=Y_QVG(\tau)$. Integrating $P\dd t=Y_QV\dd N$ gives $Q$, and $\dd\alpha/\dd\tau=G(\tau)/(N_{\max}-N_0)$, which is a function of $\alpha$ because $\alpha$ is strictly increasing in $\tau$ while $N<N_{\max}$. \hfill$\square$

\textbf{Proposition~\ref{prop:riemann}.} Let $g(t)=\mu(T(t))$. The error is $\frac1K\sum_i\int_{t_i}^{t_{i+1}}[g(t_i)-g(t)]\dd t$. (a) Each integrand is bounded by the oscillation of $g$ on $[t_i,t_{i+1}]$, which is at most the variation $V_i$ there, so the error is at most $\Delta t\sum_iV_i/K=\Delta tV/K$. (b) For a pulse of height $\Delta g=\mu(T_1)-\mu(T_0)$ starting at $s\in[t_i,t_{i+1})$ and ending at $e\in[t_j,t_{j+1})$, the start is undercounted by $(t_{i+1}-s)\Delta g$ and the end overcounted by $(t_{j+1}-e)\Delta g$; each lies in $[0,\Delta t\Delta g)$, so their difference lies in $(-\Delta t\Delta g,\Delta t\Delta g)$. (c) If the phases are independent and uniform, the two terms are independent $\mathcal U(0,\Delta t\Delta g)$ variables; their difference has mean 0 and variance $2(\Delta t\Delta g)^2/12$. \hfill$\square$

\textbf{Proposition~\ref{prop:sensor}.} With $\mu=b^2(T-T_{\min})^2$, $\partial\ln\CSL/\partial\beta=2\int(T-T_{\min})\dd t/\int(T-T_{\min})^2\dd t$, which equals $2/(\bar T-T_{\min})$ for a steady trace. For noise, $\mathrm E(T+\varepsilon-T_{\min})^2=(T-T_{\min})^2+\sigma_n^2$; the cross term $2(T-T_{\min})\varepsilon$ has mean zero and its sum over $n$ independent readings has standard deviation proportional to $\sqrt n$, so its relative contribution falls as $1/\sqrt n$. Quantisation is handled in the same way with variance $\delta_q^2/12$. \hfill$\square$

\textbf{Proposition~\ref{prop:auc}.} Take a spoiled trip $(z_1,s_1)$ and an unspoiled trip $(z_0,s_0)$. With $u=z_1$, $v=-z_0$ and $w=s_1-s_0=u+v+e_1-e_0$, $\AUC=\Pr(w>0\mid u>0,v>0)=4\Pr(u>0,v>0,w>0)$. The vector $(u,v,w)$ is zero-mean normal with $\mathrm{corr}(u,v)=0$ and $\mathrm{corr}(u,w)=\mathrm{corr}(v,w)=\sigma_z/\sqrt{2(\sigma_z^2+\sigma_\mathrm{e}^2)}=\rho$. The trivariate orthant probability is $\tfrac18+\tfrac{1}{4\pi}(\arcsin0+2\arcsin\rho)$, so $\AUC=\tfrac12+\tfrac{2}{\pi}\arcsin\rho$, which is Eq.~\eqref{eq:binorm}. \hfill$\square$


